% 04-cai-dat.tex — Chương 4: cài đặt hệ thống.
% 4.1 kiến trúc chung (viết dần), 4.2 trở đi: mỗi nhóm chức năng một mục.
\chapter{Cài đặt hệ thống}

\section{Kiến trúc tổng thể}

Phía máy chủ được xây dựng trên .NET 10 theo kiến trúc phân lớp (Clean Architecture) gồm bốn dự án:
\code{Domain} chứa thực thể và quy tắc nghiệp vụ, \code{Application} chứa các ca sử dụng dưới dạng
lệnh và bộ xử lý, \code{Infrastructure} cài đặt truy cập cơ sở dữ liệu, gửi thư, hàng đợi, và
\code{Api} nhận yêu cầu HTTP. Lớp trong không phụ thuộc lớp ngoài: \code{Application} chỉ khai báo
giao diện (ví dụ \code{IEmailSender}), \code{Infrastructure} cài đặt và được đăng ký qua cơ chế tiêm
phụ thuộc. Phía giao diện là ứng dụng React viết bằng TypeScript, dùng thư viện Ant Design.

\chuaviet[Sơ đồ triển khai, cấu trúc thư mục và các thành phần hạ tầng (SQL Server, RabbitMQ, Redis,
MinIO) sẽ được bổ sung.]

% =====================================================================================
\section{Cài đặt nhóm chức năng Tài khoản}
\label{sec:cai-dat-tai-khoan}

\subsection{Luồng xử lý một yêu cầu phía máy chủ}

Mọi API của nhóm Tài khoản đi qua cùng một chuỗi lớp. Lấy chức năng đăng ký làm ví dụ:

\begin{enumerate}
  \item \code{AuthController} nhận yêu cầu \code{POST /api/auth/register}, chuyển nội dung thành lệnh
  \code{RegisterCommand} và gửi qua thư viện MediatR. Bộ điều khiển không chứa logic nghiệp vụ.
  \item Trước khi tới bộ xử lý, \code{ValidationBehavior} chạy các luật kiểm tra viết bằng FluentValidation
  (BR01--BR05). Dữ liệu sai thì dừng ngay và trả lỗi 400 kèm mã lỗi theo từng trường.
  \item \code{RegisterCommandHandler} kiểm tra email chưa tồn tại, băm mật khẩu và gọi phương thức
  \code{User.Register} của thực thể. Thực thể tự đặt trạng thái chưa xác thực và ghi nhận sự kiện miền
  \code{UserRegistered}.
  \item Repository lưu thực thể. Ngay trước khi ghi xuống cơ sở dữ liệu, sự kiện miền được đổi thành
  thông điệp \code{SendVerificationEmailMessage} và ghi vào bảng outbox trong cùng giao dịch.
  \item Lỗi nghiệp vụ được ném dưới dạng ngoại lệ có mã; một middleware chung đổi chúng thành phản hồi
  ProblemDetails với mã HTTP và trường \code{code} tương ứng (Bảng~\ref{tab:tk-ma-loi}).
\end{enumerate}

Cách tổ chức này giữ mỗi ca sử dụng trong một thư mục riêng (lệnh, luật kiểm tra, bộ xử lý, kết quả),
dễ tìm và dễ kiểm thử độc lập. Truy vấn dữ liệu được viết bằng mẫu Specification (thư viện
Ardalis.Specification), ví dụ \code{UserByEmailSpec}, để điều kiện truy vấn có tên và dùng lại được.

\subsection{Gửi email bất đồng bộ qua outbox}

Gửi thư là thao tác chậm và có thể lỗi, nên không được thực hiện trong yêu cầu đăng ký. Hệ thống dùng
mẫu Transactional Outbox~\cite{outbox} do thư viện MassTransit cung cấp:

\begin{itemize}
  \item \textbf{Phía gửi:} thông điệp được ghi vào bảng \code{OutboxMessage} cùng giao dịch với tài
  khoản mới. Hoặc cả hai cùng được lưu, hoặc không có gì được lưu; không xảy ra trường hợp tài khoản đã
  tạo mà thông điệp bị mất. Một tiến trình nền đọc bảng outbox và đẩy thông điệp lên RabbitMQ.
  \item \textbf{Phía nhận:} \code{SendVerificationEmailConsumer} nhận thông điệp và chuyển cho bộ xử lý
  \code{SendVerificationEmailCommandHandler}. Bảng inbox ghi nhận thông điệp đã xử lý nên mỗi thông
  điệp chỉ được xử lý một lần; lỗi thì thử lại sau 1, 5 và 15 giây.
\end{itemize}

Bộ xử lý gửi thư thực hiện năm bước: đọc lại tài khoản theo mã và bỏ qua nếu tài khoản không còn ở
trạng thái chưa xác thực; bỏ qua nếu vừa gửi thư cùng loại trong 60 giây (BR08); đánh dấu các token cũ
còn hiệu lực là đã dùng; sinh token mới gồm 32 byte ngẫu nhiên (mã hóa Base64Url), lưu bản băm SHA-256
kèm thời hạn; cuối cùng gửi thư bằng thư viện MailKit, liên kết trong thư có dạng
\code{<địa chỉ giao diện>/verify-email?token=<token gốc>}.

Thông điệp chỉ mang mã người dùng, không mang token (BR07). Thông điệp nằm ở nhiều nơi ngoài tầm
kiểm soát của ứng dụng: bảng outbox, hàng đợi RabbitMQ và hàng đợi lỗi (giữ lại vĩnh viễn nếu xử lý
thất bại hết số lần thử). Nếu token nằm trong thông điệp, ai đọc được các nơi đó sẽ có liên kết xác
thực hoặc đặt lại mật khẩu còn hạn. Sinh token ngay lúc gửi thư khiến chuỗi token gốc chỉ tồn tại trong
email; cơ sở dữ liệu chỉ có bản băm, giống cách lưu mật khẩu.

\subsection{Phiên đăng nhập: access token và refresh token}

Sau khi đăng nhập thành công, hệ thống cấp hai loại token với vai trò khác nhau:

\begin{itemize}
  \item \textbf{Access token} là JWT~\cite{rfc7519} ký bằng HMAC-SHA256, hiệu lực 15 phút, chứa mã người
  dùng, email, họ tên, vai trò và danh sách quyền. Giao diện gửi token này trong tiêu đề
  \code{Authorization} của mỗi yêu cầu. Máy chủ kiểm tra chữ ký và thời hạn mà không cần truy vấn cơ sở
  dữ liệu.
  \item \textbf{Refresh token} là chuỗi ngẫu nhiên, hiệu lực 7 ngày, nằm trong cookie
  \code{confhub\_refresh} với các thuộc tính \code{HttpOnly} (mã JavaScript không đọc được),
  \code{Secure}, \code{SameSite=Strict} và \code{Path=/api/auth} (chỉ gửi kèm các API tài khoản).
  Cơ sở dữ liệu chỉ lưu bản băm của token này.
\end{itemize}

Mỗi lần làm mới, refresh token cũ bị đánh dấu đã dùng và một token mới được cấp (xoay vòng, BR12).
Nếu một token đã dùng lại được gửi lên, nhiều khả năng token đã bị đánh cắp và đang được hai bên cùng
sử dụng; bộ xử lý \code{RefreshSessionCommandHandler} khi đó thu hồi toàn bộ refresh token của người
dùng, buộc mọi thiết bị đăng nhập lại~\cite{owasp-session}.

Phân quyền dựa trên danh sách quyền trong access token. Các API cần quyền được đánh dấu bằng thuộc
tính \code{[MustHavePermission]}; một bộ cung cấp chính sách tự sinh chính sách kiểm tra quyền tương
ứng. Mặc định mọi API đều yêu cầu đăng nhập; chỉ các API được đánh dấu riêng (đăng ký, đăng nhập, làm
mới phiên, quên mật khẩu\ldots) mới cho phép truy cập khi chưa đăng nhập.

\subsection{Các biện pháp bảo mật}

\begin{itemize}
  \item \textbf{Mật khẩu} được băm bằng \code{PasswordHasher} của ASP.NET Core (thuật toán PBKDF2 có
  muối)~\cite{owasp-password}. Khi email không tồn tại, hệ thống vẫn thực hiện một lần băm giả để thời
  gian phản hồi không để lộ email nào có tài khoản.
  \item \textbf{Token trong email và refresh token} chỉ lưu bản băm SHA-256. Vì token là chuỗi ngẫu nhiên
  256 bit nên không cần thuật toán băm chậm như mật khẩu.
  \item \textbf{Không lộ thông tin tài khoản:} cùng thông báo cho sai email và sai mật khẩu; cùng phản
  hồi cho gửi lại email và quên mật khẩu (BR09, BR10).
  \item \textbf{Giới hạn tần suất} theo địa chỉ IP bằng bộ giới hạn có sẵn của ASP.NET Core, cửa sổ cố
  định một phút (BR15); vượt giới hạn trả mã 429.
  \item \textbf{Access token chỉ nằm trong bộ nhớ} của trang (không lưu \code{localStorage}), giảm
  thiệt hại nếu trang bị chèn mã độc.
\end{itemize}

\subsection{Phía giao diện}

Mã giao diện được chia theo nhóm chức năng (\code{features/auth}); trong mỗi nhóm, các lớp chỉ gọi lớp
ngay dưới: trang (page) $\rightarrow$ thành phần giao diện (component) $\rightarrow$ hook
$\rightarrow$ hàm gọi API $\rightarrow$ \code{apiClient}. Các hook dùng thư viện TanStack Query để quản
lý trạng thái đang tải, lỗi và kết quả của mỗi lần gọi API.

\textbf{Lớp gọi API dùng chung.} \code{apiClient} là một đối tượng axios với hai bộ chặn
(interceptor). Bộ chặn yêu cầu tự gắn access token; bộ chặn phản hồi đổi mọi lỗi thành một kiểu thống
nhất gồm mã HTTP, mã lỗi nghiệp vụ và lỗi theo từng trường. Khi một API trả 401 vì access token hết
hạn, bộ chặn tự gọi làm mới phiên một lần rồi gửi lại yêu cầu; làm mới thất bại thì xóa phiên và đưa
người dùng về trang đăng nhập kèm thông báo phiên đã hết hạn.

\textbf{Khôi phục phiên khi tải lại trang.} Vì access token chỉ nằm trong bộ nhớ, tải lại trang sẽ làm
mất token. Khi ứng dụng khởi động, thành phần \code{SessionBootstrap} hiển thị vòng chờ và gọi API làm
mới phiên; trình duyệt tự gửi cookie refresh. Còn phiên thì nhận access token mới, không còn thì người
dùng là khách. Trang chỉ được vẽ sau bước này, nên thanh tiêu đề hiển thị đúng trạng thái ngay từ đầu.

\textbf{Chỉ một yêu cầu làm mới tại một thời điểm.} Do refresh token xoay vòng, hai yêu cầu làm mới gửi
cùng lúc bằng cùng một cookie sẽ bị máy chủ coi là token bị đánh cắp (BR12). Tình huống này xảy ra khi
nhiều API cùng nhận 401, hoặc khi React chạy hiệu ứng hai lần ở chế độ phát triển. Hàm
\code{refreshSession} giữ lời gọi đang chạy trong một biến dùng chung: lời gọi thứ hai nhận lại đúng
lời gọi đang chạy thay vì gửi yêu cầu mới. Ngoài ra, mỗi lần phiên thay đổi có một số đếm tăng lên;
kết quả làm mới trả về muộn, sau khi người dùng đã đăng xuất hoặc đăng nhập tài khoản khác, sẽ bị bỏ qua.

\textbf{Song ngữ và thông báo lỗi.} Toàn bộ chữ hiển thị nằm trong hai từ điển tiếng Việt và tiếng
Anh. Mã lỗi máy chủ trả về (ví dụ \code{EmailTaken}) được dùng trực tiếp làm khóa tra từ điển, nên máy
chủ không cần biết người dùng đang xem ngôn ngữ nào. Lỗi theo từng trường được gắn vào đúng ô nhập; khi
người dùng đổi ngôn ngữ, các thông báo lỗi đang hiển thị cũng được dịch lại.

\subsection{Giao diện chức năng}

Các màn hình của nhóm Tài khoản nằm trong khung giao diện công khai, dạng thẻ ở giữa trang. Các hình
dưới đây được chụp trên môi trường phát triển với máy chủ thật; thư được xem qua công cụ smtp4dev.

\begin{figure}[H]
  \centering
  \includegraphics[width=0.48\textwidth]{screenshots/hinh-02-dang-ky-kiem-tra-du-lieu.png}
  \caption{Màn hình đăng ký với vai trò Ban tổ chức}
  \label{fig:gd-dang-ky}
\end{figure}

Hình~\ref{fig:gd-dang-ky} là màn hình đăng ký khi người dùng chọn vai trò Ban tổ chức và bấm ``Đăng
ký'' mà chưa nhập thông tin. Ô ``Đơn vị / công ty'' và ghi chú cần quản trị viên duyệt chỉ xuất hiện với
Ban tổ chức và Nhà cung cấp. Lỗi được báo ngay dưới từng ô; dưới ô mật khẩu có dòng gợi ý quy tắc mật
khẩu (BR02).

\begin{figure}[H]
  \centering
  \includegraphics[width=0.62\textwidth]{screenshots/hinh-03-kiem-tra-email.png}
  \caption{Màn hình kiểm tra email sau khi đăng ký}
  \label{fig:gd-kiem-tra-email}
\end{figure}

Sau khi đăng ký thành công, người dùng được chuyển sang màn hình ở Hình~\ref{fig:gd-kiem-tra-email}.
Màn hình hiển thị địa chỉ vừa đăng ký và nút gửi lại thư; nút bị khóa 60 giây kể từ khi vào trang và
hiển thị số giây còn lại (BR08). Liên kết ``Đăng ký lại'' dành cho trường hợp nhập sai email.

\begin{figure}[H]
  \centering
  \begin{minipage}[t]{0.48\textwidth}
    \centering
    \includegraphics[width=\textwidth]{screenshots/hinh-04-xac-thuc-thanh-cong.png}
  \end{minipage}\hfill
  \begin{minipage}[t]{0.48\textwidth}
    \centering
    \includegraphics[width=\textwidth]{screenshots/hinh-05-xac-thuc-lien-ket-hong.png}
  \end{minipage}
  \caption{Kết quả xác thực email: thành công (trái) và liên kết đã dùng (phải)}
  \label{fig:gd-xac-thuc}
\end{figure}

Hình~\ref{fig:gd-xac-thuc} là hai kết quả của trang xác thực email. Khi người dùng bấm liên kết trong
thư, trang tự gửi token lên máy chủ và hiển thị kết quả máy chủ trả về. Bấm lại cùng liên kết lần thứ
hai thì token đã được dùng, trang báo liên kết không còn hiệu lực và cho nhập email để nhận liên kết
mới. Với tài khoản Ban tổ chức hoặc Nhà cung cấp, trang báo đã xác thực và đang chờ quản trị viên duyệt.

\begin{figure}[H]
  \centering
  \begin{minipage}[t]{0.48\textwidth}
    \centering
    \includegraphics[width=\textwidth]{screenshots/hinh-06-dang-nhap-sai-thong-tin.png}
  \end{minipage}\hfill
  \begin{minipage}[t]{0.48\textwidth}
    \centering
    \includegraphics[width=\textwidth]{screenshots/hinh-07-dang-nhap-tieng-anh.png}
  \end{minipage}
  \caption{Màn hình đăng nhập báo sai thông tin, tiếng Việt (trái) và tiếng Anh (phải)}
  \label{fig:gd-dang-nhap}
\end{figure}

Hình~\ref{fig:gd-dang-nhap} là màn hình đăng nhập khi nhập sai mật khẩu. Thông báo không cho biết sai
email hay sai mật khẩu (BR10). Khi người dùng chuyển sang tiếng Anh, toàn bộ chữ trên trang, kể cả
thông báo lỗi do máy chủ trả về, được đổi ngôn ngữ mà không cần gửi lại yêu cầu.

\begin{figure}[H]
  \centering
  \includegraphics[width=0.75\textwidth]{screenshots/hinh-01-header-menu-tai-khoan.png}
  \caption{Thanh tiêu đề sau khi đăng nhập, mở menu tài khoản}
  \label{fig:gd-menu}
\end{figure}

Sau khi đăng nhập, thanh tiêu đề thay hai nút ``Đăng nhập'', ``Đăng ký'' bằng tên người dùng
(Hình~\ref{fig:gd-menu}). Bấm vào tên sẽ mở menu gồm email, vai trò và mục ``Đăng xuất''. Tải lại trang
vẫn giữ trạng thái đăng nhập nhờ cơ chế khôi phục phiên đã trình bày ở trên.

\begin{figure}[H]
  \centering
  \begin{minipage}[t]{0.48\textwidth}
    \centering
    \includegraphics[width=\textwidth]{screenshots/hinh-10-dat-lai-mat-khau-thanh-cong.png}
  \end{minipage}\hfill
  \begin{minipage}[t]{0.48\textwidth}
    \centering
    \includegraphics[width=\textwidth]{screenshots/hinh-11-dang-nhap-nen-toi.png}
  \end{minipage}
  \caption{Đặt lại mật khẩu thành công (trái) và màn hình đăng nhập ở nền tối (phải)}
  \label{fig:gd-khac}
\end{figure}

Hình~\ref{fig:gd-khac} bên trái là kết quả sau khi đặt mật khẩu mới từ liên kết trong thư; mọi phiên
đăng nhập cũ của tài khoản đã bị thu hồi (BR14). Bên phải là giao diện ở nền tối: người dùng chọn
nền sáng hoặc tối trên thanh tiêu đề, lựa chọn được lưu lại cho lần mở sau.

\subsection{Kiểm thử}

Mỗi ca sử dụng được kiểm thử tự động ở cả hai phía. Phía máy chủ có kiểm thử đơn vị cho bộ xử lý,
luật kiểm tra dữ liệu, và kiểm thử tích hợp gọi trực tiếp các API, trong đó có các trường hợp: token hết hạn, token đã dùng, xoay vòng và phát hiện refresh token bị dùng lại, giới
hạn tần suất. Phía giao diện kiểm thử lớp gọi API (đổi lỗi, làm mới phiên một lần, chỉ một yêu cầu làm
mới) và các luồng màn hình chạy trên bộ định tuyến thật với máy chủ giả lập.

\chuaviet[Bảng tổng hợp số trường hợp kiểm thử và kết quả sẽ được bổ sung khi hoàn thiện các nhóm chức năng.]
